\documentclass[10pt,a4paper]{amsart}
\usepackage[margin=19mm]{geometry}
\usepackage{amsmath,amssymb,mathtools}
\usepackage[scr=rsfs,cal=euler]{mathalfa}
\usepackage{xfrac}
\usepackage[unicode,hidelinks,bookmarks=false]{hyperref}
\newcommand{\ie}{i.e.\ }
\newcommand{\cf}{cf.\ }
\newcommand{\resp}{respectively\ }
\newcommand{\Subsection}{\subsection}
\providecommand{\nobreakdash}{\nobreak\hbox{-}}
\setlength{\emergencystretch}{2em}
\multlinegap0pt
\usepackage[normalem]{ulem}
\usepackage{xcolor}
\usepackage{amssymb}
\usepackage{float}
\usepackage{enumerate}
\newtheorem{theorem}{Theorem}
\newtheorem{lemma}{Lemma}
\newcommand{\bb}{\mathbb}
\newcommand{\irr}{\mathrm{Irr}}
\newcommand{\mbf}{\mathbf}
\newcommand{\mcal}{\mathcal}
\newcommand{\mrm}{\mathrm}
\newcommand{\ov}{\overline}
\title{Perfect and Multiple state Transfer in Oriented Cayley Graphs}
\author{Ada Chan, Venkata Raghu Tej Pantangi, Andriaherimanana Sarobidy Razafimahatratra, Peter Sin}
\date{Source: arXiv:2510.09873v2 (CC BY 4.0)}
\begin{document}
\maketitle

\noindent\emph{Source and licence.} Perfect and Multiple state Transfer in Oriented Cayley Graphs. Version arXiv:2510.09873v2 (CC BY 4.0), distributed by arXiv under the Creative Commons Attribution 4.0 International licence, http://creativecommons.org/licenses/by/4.0/, as recorded in the arXiv licence metadata for this exact version and shown on the abstract page https://arxiv.org/abs/2510.09873v2. Full licence notice: \texttt{licenses/arxiv-2510.09873-CC-BY-4.0.txt}.\par\smallskip

\noindent\textit{Editorial scope.} Counted unit: the article's theorem thm:wreath (source0.tex lines 819--829). The article's complete original proof of it is delivered with it (source0.tex lines 830--888, one \texttt{proof} environment of 59 lines). No mathematics is written, repaired or re-derived: the statement and the proof are the article's own lines, in the article's own notation, and no retained line is substituted or re-flowed.  Retained uncounted as the source-local context the proof needs: lines 371--379; lines 764--766; lines 771--776; lines 795--807.  Nothing was omitted from any retained span, so the package carries no omission record.  No retained line contains a citation, so the wrapper carries no bibliography and no bibliography entry.  No retained line contains a figure environment, an image-inclusion command or drawing code, so no image is delivered and the package declares no asset dependency.  Editorial formatting only, in addition: an AMS \texttt{amsart} preamble that restates the article's own declarations and macros used by the retained lines, namely the environments \texttt{theorem}, \texttt{lemma} on separate counters, together with the standard packages those lines need.  The retained environments print in this document as Theorem 1, Lemma 1 in order of appearance, whereas the article prints the same items with the numbering of its own full document.  The complete native source of the article as distributed in the arXiv e-print for this exact version is bundled unchanged as \texttt{sources/186840/source0.tex}.
\begin{theorem}\label{thm:characterization}
PST occurs  from $e$ to $z$ at time $\tau$ in an oriented normal Cayley graph $\mrm{Cay}(G, C)$ if and only if the following hold.
\begin{enumerate}
\item $z\in Z(G)$.
\item \label{Cond:chi}
For $\chi \in \irr(G)$,
\[\dfrac{\chi(z)}{\chi(e)}=  \exp\left({\tau\ \dfrac{\chi(C)-\ov{\chi(C)}}{\chi(e)}}\right).\]
\end{enumerate}
\end{theorem}  

\begin{equation}\label{eq:charwreathproduct}
\psi_{f,\Lambda}:= \mrm{Ind}_{G \wr S_{\mcal{B}_{f}}}^{G \wr S_{n}} \left[ \bigotimes_{\phi \in \irr(G)} \left(\widetilde{\phi^{|f^{-1}(\phi)|}_{\Lambda(\phi)}} \right)\right] \in \irr(G\wr S_{n}).
\end{equation}

\begin{lemma}\label{lem:centralcharacterc1}
Consider $f:[n]\to \irr(G)$ and a partition-valued function $\Lambda$ on $\irr(G)$ satisfying $|\Lambda(\phi)|=|f^{-1}(\phi)|$, for all $\phi \in \irr(G)$. Then, we have
\begin{equation*}
\dfrac{\psi_{f,\Lambda}(((x,e,\ldots, e); (1)))-\ov{\psi_{f,\Lambda}(((x,e,\ldots, e); (1)))}}{\psi_{f,\Lambda}(1)} = \sum\limits_{\mu \in \irr(G)}  \dfrac{|f^{-1}(\mu)| (\mu(x)-\ov{\mu(x)})}{n\mu(1)}.
\end{equation*}
\end{lemma}

\begin{lemma}\label{lem:centralcharacterc2}
Let $f:[n]\to \irr(G)$ be a function, let  $\Lambda$ a partition-valued function on $\irr(G)$ satisfying $|\Lambda(\phi)|=|f^{-1}(\phi)|$, for all $\phi \in \irr(G)$, and let $\pi \in S_{n}$ be an $n$-cycle.
Then the following statements are true for all $\mbf{x}\in G^{n}$.
\begin{enumerate}
    \item If $|f([n])|>1$, then 
    \[\dfrac{\psi_{f,\Lambda}((\mbf{x}; \pi)-\ov{\psi_{f,\Lambda}((\mbf{x}; \pi)}}{\psi_{f,\Lambda}(1)}=0;\] and
    \item if $f([n])=\{\phi\}$, then 
    \[
    |G|^{n-1} \times\dfrac{\psi_{f,\Lambda}((\mbf{x}; \pi)-\ov{\psi_{f,\Lambda}((\mbf{x}; \pi)}}{\psi_{f,\Lambda}(1)} = k_{f,\Lambda}|Z(G)|^{(n-1)} \dfrac{\phi\left( \mbf{x}_{\pi}  \right)-\ov{\phi\left( \mbf{x}_{\pi}  \right)}}{\phi(1) \times (n-1)!},
    \]
    for some $k_{f,\Lambda} \in \bb{Z}$. 
\end{enumerate}
\end{lemma}

\begin{theorem}\label{thm:wreath}
Let $n\geq 2$ be a natural number and let $X:=\mrm{Cay}(G,C)$ be an oriented normal Cayley graph that admits PST from $e$ to $z\in Z(G)$, at a time $\tau$. 
Let 
\begin{subequations}\label{eq:wreathconnex}
\begin{align}
\tilde{C}_{1} &=  \bigcup\limits_{i \in [n]} \left\{((x_{1}, \ldots x_{n}); (1))\ :\ x_{i}\in C\ \& \ x_{j}=e,\ \forall j\neq i\right\}\\
\tilde{C}_{2} &= \left\{(\mbf{x};\pi)\ :\ \text{$\pi$ is an $n$-cycle}\ \&\ \mbf{x}_{\pi}\in C\right\}.
\end{align} 
\end{subequations}
Then $X \wr n := \mrm{Cay}(G\wr S_{n},\ \tilde{C}_{1}\cup \tilde{C}_{2})$ admits PST at time $\tau$ from $(e,e,\ldots, e)$ to $(z,z,\ldots, z)$.
\end{theorem}

\begin{proof}
By Theorem~\ref{thm:characterization}, it suffices to show that 
\[\dfrac{\ov{\psi(z,\ldots, z ) }}{\psi(1)} = \exp\left(\tau \theta_{\psi}\right),\] for all $\psi \in \irr(G \wr S_{n})$. Here,
 \begin{equation}\label{eq:tpsiwr}
\theta_{\psi} = \dfrac{\psi(\tilde{C}_{1})+\psi(\tilde{C}_{2})-\ov{\psi(\tilde{C}_{1})-\psi(\tilde{C}_{2})}}{\psi(1)}.
\end{equation}

Consider some $\psi \in \irr(G \wr S_{n})$.
Following the description of $\irr(G \wr S_{n})$ given at the start of this section, given $\psi \in \irr(G\wr S_{n})$, there exists $f:[n] \to \irr{G}$ and a partition-valued function $\Lambda$ on $\irr(G)$ with $|\Lambda(\phi)|=|f^{-1}(\phi)|$ (for all $\phi \in \irr(G)$) such that $\psi=\psi_{f,\Lambda}$ (as defined in \eqref{eq:charwreathproduct}). We first compute $\theta_{\psi_{f,\Lambda}}$.

\paragraph{Case 1:} Assume that $|f([n])|>1$. Using Lemma~\ref{lem:centralcharacterc2}, it follows that
\[\theta_{\psi_{f, \Lambda}} = \dfrac{\psi_{f, \Lambda}(\tilde{C}_{1}))-\ov{\psi_{f, \Lambda}(\tilde{C}_{1})}}{\psi_{f, \Lambda}(1)}.\] We note that every element of $\tilde{C}_{1}$ is conjugate to an element of the form $((x,e,\ldots ,e) ; (1))$ for some $x\in C$, and that the number of $S_{n}$-conjugates in $\tilde{C}_{1}$ of the element $((x,e,\ldots ,e) ; (1))$ (with $x \in C$) is $n$ . It now follows from Lemma~\ref{lem:centralcharacterc1} that 
\begin{align*}
\dfrac{\psi_{f, \Lambda}(\tilde{C}_{1})}{\psi_{f, \Lambda}(1)} = n\sum\limits_{\phi \in \irr(G)}  \dfrac{|f^{-1}(\phi)|\phi(C)}{n\phi(1)},
\end{align*}
and therefore, if $\theta_{\phi}$ is the eigenvalue associated with $\phi$, of $\mrm{Cay}(G,C)$, we have
\begin{equation}\label{eq:c1sum}
\theta_{\psi_{f, \Lambda}} = \sum\limits_{\phi \in \irr(G)} |f^{-1}(\phi)| \theta_{\phi},\ \text{provided $|f([n]|>1$.}
\end{equation}

 \paragraph{Case 2:} Now, we assume that $f([n])=\{\phi\}$. Set $\lambda:=\Lambda(\phi)$. Fix an $n$-cycle, $\sigma \in S_{n}$. We note that every element in $\tilde{C}_{2}$ is conjugate to an element of the form $(\mbf{x}; \sigma)$ with $\mbf{x}_{\sigma} \in C$. We count that for each $c \in C$, that $X_{c}:=\{(\mbf{x}; \sigma) \ : \mbf{x}_{\sigma} = c \}$ is a set of size $|G|^{n-1}$.  Since every element of $X_{c}$ lies in the same conjugacy class, using Lemma~\ref{lem:centralcharacterc2}, we have
 \begin{align}\label{eq:formulac2}
     \sum\limits_{g \in X_{c}}\dfrac{\psi_{f,\Lambda}(g)-\ov{\psi_{f,\Lambda}(g)}}{\psi_{f, \Lambda}(1)} & =k_{f, \lambda} |Z(G)|^{n-1} \dfrac{\phi(c)-\ov{\phi(c)}}{\phi(1)\times (n-1)!},
 \end{align}
 for some $k_{f,\Lambda}\in \bb{Z}$.
 We recall that there are exactly $(n-1)!$ number of $n$-cycles in $S_{n}$. Thus, we have
\begin{align*}
\dfrac{\psi_{f, \Lambda}(\tilde{C}_{2})-\ov{\psi_{f, \Lambda}(\tilde{C}_{2})}}{\psi_{f, \Lambda}(1)} &= (n-1)!\sum\limits_{c \in C}\sum\limits_{g \in X_{c}}\dfrac{\psi_{f,\Lambda}(g)-\ov{\psi_{f,\Lambda}(g)}}{\psi_{f, \Lambda}(1)} \\
&= \sum\limits_{c \in C}k_{f,\Lambda}|Z(G)|^{n-1} \dfrac{\phi(c) - \ov{\phi(c)}}{\phi(1)} \ \text{(using \eqref{eq:formulac2})}\\
&=k_{f,\Lambda}|Z(G)|^{n-1}\theta_{\phi}.
\end{align*}
Arguing as in Case 1, using Lemma~\ref{lem:centralcharacterc1}, we have
\[\dfrac{\psi_{f, \Lambda}(\tilde{C}_{1})-\ov{\psi_{f, \Lambda}(\tilde{C}_{1})}}{\psi_{f, \Lambda}(1)}=n\theta_{\phi},\] and therefore, we have 

\begin{equation}\label{eq:c2sum}
\theta_{\psi_{f, \Lambda}} = n \theta_{\phi} + k_{f,\Lambda}|Z(G)|^{n-1}\theta_{\phi},\ \text{, for some $k_{f,\Lambda}\in \bb{Z}$, provided $f([n])=\{\phi\}$.}
\end{equation}

Using, \eqref{eq:tpsiwr}, \eqref{eq:c1sum} and \eqref{eq:c2sum}, we have 
\begin{align}\label{eq:taupsi}
\exp\left( \tau \theta_{\psi_{f,\Lambda}}\right) &= \begin{cases} \prod\limits_{\phi \in \irr(G)} \left[ \exp\left( \tau \theta_{\phi} \right) \right]^{|f^{-1}(\phi)|} & \text{if $|f([n])|>1$,}\\
\exp{(\tau\theta_{\phi})}^{n} \times \exp(\tau \theta_{\phi})^{k_{\psi} |Z(G)|^{n-1}} &\text{if $f([n])=\{\phi\}$.} 
\end{cases}
\end{align}
Since $\mrm{Cay}(G, C)$ admits PST from $e$ to $z$ at time $\tau$, by Theorem~\ref{thm:characterization}, we have \[\exp\left( \tau \theta_{\phi} \right) = \frac{\ov{\phi(z)}}{\phi(1)},\] for all $\phi \in \irr(G)$. Applying the above in \eqref{eq:taupsi}, yields

\begin{align*}
\exp\left( \tau \theta_{\psi_{f,\Lambda}}\right) & = \begin{cases}
\prod\limits_{\phi \in \irr(G)} \left( \frac{\ov{\phi(z)}}{\phi(1)} \right)^{|f^{-1}(\phi)|}  & \text{if $|f([n])|>1$,} \\
\left( \dfrac{\ov{\phi(z)}}{\phi(1)} \right)^{n} \times \left(\dfrac{\ov{\phi(z)}}{\phi(1)}\right)^{k_{f,\Lambda}|Z(G)|^{n-1}} & \text{if $f([n])=\{\phi\}$.}
\end{cases}
\end{align*} 
Since $z \in Z(G)$ and since $\dfrac{\ov{\phi(z)}}{\phi(1)}$ is a $|z|$-th root of unity, it follows that $\dfrac{\ov{\phi(z)}}{\phi(1)}^{|Z(G)|}=1$. Therefore, we have
\begin{align*}
\exp\left( \tau \theta_{\psi_{f,\Lambda}}\right) & = \prod\limits_{\phi \in \irr(G)} \left( \dfrac{\ov{\phi(z)}}{\phi(1)} \right)^{|f^{-1}(\phi)|} \\
&= \dfrac{\ov{\psi_{f,\Lambda}(z,\ldots, z)} }{\psi_{f,\Lambda}(1)}.
\end{align*} 
The result now follows from Theorem~\ref{thm:characterization}.
\end{proof}

\end{document}
